\bookpreface{Preface}

This book grew out of the way I used to teach ML. I came to machine learning from physics, so I think about it quite differently from the average computer scientist or statistician. I haven't seen another book approach the subject quite this way.

There's nothing particularly new here apart from the route I take to get to the results. A lot of that comes from my training in physics, especially theoretical physics. It also has a lot to do with how much I love \emph{The Feynman Lectures on Physics}. Those lectures changed the way I think about physics as a whole.

I didn't write this book for someone who's just starting out. I think you'd have a hard time learning the subject from it, because I've left out a lot of the classic material and much of the structure that helps you through a first course. But if this is your second or third pass, you may find another way to look at familiar ideas and understand \emph{how} they're used in practice.

You can read the first few chapters alongside \href{https://amlbook.com/}{\emph{Learning from Data}}, by Yaser S. Abu-Mostafa, Malik Magdon-Ismail, and Hsuan-Tien Lin. There you'll find a more traditional course that teaches the mathematical foundations of the field.

I'm especially fond of the distinction I draw between \emph{generative learning} and \emph{discriminative learning}. Most books mention both approaches, but I think a lot of readers come away without much of a feel for what it means to take one path or the other.
